% Ad Atticum 13.29 --- headnote
% ad-atticum-13-29, 27 May 45 BC, Tusculanum

\textit{Cicero to Atticus, written at the Tusculan villa on 27 May 45
BC --- Perseus dateline \textit{Scr. in Tusculano vi K. Iun. a. 709
(45)}. One in the daily series of late-May letters from Tusculum
that revolve around two intertwined projects: the revision and
Varro-dedication of the \textit{Academica}, and the increasingly
elaborate financial machinery for buying the suburban gardens
across the Tiber where Cicero means to build a shrine
(\textit{fanum}) to his dead daughter Tullia. Here the gardens
question dominates: Chrysippus has been out to inspect the
Scapulan property, the villa is what Cicero expected, the
bathrooms are tolerable, a covered walk will have to be added, and
the right grove for the shrine has, awkwardly, become a busy
spot.}

\textit{Two Greek phrases punctuate the page. The first,
\textit{aphidruma} ``shrine,'' is the technical term Cicero keeps
reaching for when he describes what the gardens are for: not a
pleasure-house but a sanctuary. The second is the Homeric
\textit{ton typhon mou pros theōn tropophorēson} (after
\textit{Od.} or \textit{Il.}, with \textit{tropophoreō} from
Deut. 1.31 in the Septuagint --- but Cicero is reaching for the
Greek verb of patient bearing): ``in the gods' name, bear with
my vanity.'' He is asking Atticus, in Greek, not to mock him for
caring so much about the symbolic prestige (\textit{typhos}) of
the location. The financing turns on Faberius's outstanding debt:
if Atticus can ``clear that debt,'' Cicero will outbid Otho for
the Scapulan gardens; if not, he falls back on Clodia's, where
the Dolabella debt may suffice for cash on the spot. The closing
exchange about ``Cicero's letter'' refers to young Marcus's
complaints from his Athenian studies. The crux \textit{misissem}
I mark with daggers, following the editorial tradition.}
